\subsection{AI Systems Learning from Moral Disagreements and Conflicts}
\label{sec:5.3.3}
Section~\ref{sec:2.1.1} covers what an AI system does when its own values conflict with each other — a hierarchy among them, a case base of past resolutions, an audit process for catching drift between principle and behavior. A different problem sits next to that one: what a system does when its judgment conflicts with someone else's, human or artificial, and neither party is simply wrong. Learning from that kind of disagreement, instead of resolving it by fiat toward whichever party holds authority, is its own design target.

IBM's Project Debater argued live, unscripted policy questions against human champions, drawing on a large corpus of newspaper and journal articles to construct and rebut arguments in real time \autocite{slonim2021autonomous}. It is the clearest existing demonstration that a system can extract structured argument from open moral and political disagreement rather than retrieve a matching answer. A system built to learn from debate this way — summarizing each side, naming the ethical framework doing the work underneath it, weighing argument strength rather than detecting who won on rhetoric — builds a capacity that resolving one dilemma at a time does not: recognizing which value is in tension with which, across surface disagreements that share a structure. Simulation supplies a second route, placing agents with genuinely different objectives in a negotiation or a resource-allocation problem, which generates the kind of experience with competing interests that a system trained only on a curated distribution of correct answers never gets.

Direct feedback — a user correcting a teaching assistant's guidance, a patient's family disagreeing with a care recommendation — is the least structured of these sources and the hardest to learn from safely, because feedback from any one person is a partial, situated view and not a settled answer. A loop that updates the system's weights every time a user pushes back folds in whichever bias its most vocal users bring. One that logs the disagreement, checks whether it recurs across many different users, and only then treats it as signal is doing what that ethical auditing asks for, applied to conflict rather than to drift. A system that treats disagreement as something to smooth over risks mistaking consensus among the people it happened to sample for consensus among the people it will actually serve.
